% Part:sets-functions-relations
% Chapter: size-of-sets
% Section: reduction-alt

\documentclass[../../../include/open-logic-section]{subfiles}

\begin{document}

\olfileid{sfr}{siz}{red-alt}

\olsection{न्यूनीकरण}

\begin{editorial}
  या विभागात \olref[nen-alt]{sec} मधील फलितांशी जुळणारी अगणनीयता
  न्यूनीकरणाने सिद्ध केली आहे. \olref[nen]{sec} मधील फलितांशी जुळणारी
  किंचित अधिक विस्तृत पर्यायी आवृत्ती \olref[red]{sec} मध्ये दिली आहे.
\end{editorial}

कर्णीकरणाच्या युक्तिवादाने $\Bin^\omega$ हा !!{nonenumerable} आहे,
हे आपण सिद्ध केले. तसाच कर्णीकरणाचा युक्तिवाद वापरून $\Pow{\Nat}$ हा
!!{nonenumerable} आहे हेही दाखवले. पण $\Pow{\Nat}$ हा
!!{nonenumerable} असल्याचे सिद्ध करण्याचा आणखी एक मार्ग आहे:
\emph{$\Pow{\Nat}$ हा !!{enumerable} असेल, तर $\Bin^\omega$ हाही
!!{enumerable} असतो} हे दाखवा. $\Bin^\omega$ हा !!{nonenumerable}
आहे हे आपल्याला माहीत असल्यामुळे $\Pow{\Nat}$ हाही तसाच आहे असा
निष्कर्ष निघेल.

याला एक समस्या दुसऱ्या समस्येकडे \emph{न्यूनीत करणे} म्हणतात. या
बाबतीत $\Bin^\omega$ चे प्रगणन करण्याची समस्या आपण $\Pow{\Nat}$ चे
प्रगणन करण्याच्या समस्येकडे न्यूनीत करतो. नंतरच्या समस्येचे उत्तर---
$\Pow{\Nat}$ चे प्रगणन---पहिल्या समस्येचे उत्तर---$\Bin^\omega$ चे
प्रगणन---देईल.

संच~$B$ च्या प्रगणनाची समस्या संच~$A$ च्या प्रगणनाच्या समस्येकडे
न्यूनीत करण्यासाठी~$A$ चे प्रगणन~$B$ च्या प्रगणनात रूपांतरित करण्याची
पद्धत आपण देतो. हे करण्याचा सर्वांत सोपा मार्ग म्हणजे
$f\colon A \to B$ असे !!a{surjection} परिभाषित करणे. $x_1$, $x_2$,
\dots{} हे~$A$ चे प्रगणन असेल, तर $f(x_1)$, $f(x_2)$, \dots{} हे~$B$
चे प्रगणन करील. आपल्या बाबतीत आपण
$f\colon \Pow{\Nat} \to \Bin^\omega$ असे !!a{surjection} शोधत आहोत.

\begin{prob}
$g\colon B \to A$ असे !!{injective} फलन असेल आणि $B$ हा
!!{nonenumerable} असेल, तर $A$ हाही तसाच असतो, हे दाखवा. $g$ वापरून
~$A$ चे प्रगणन~$B$ च्या प्रगणनात कसे रूपांतरित करता येते हे दाखवून
सिद्धता द्या.
\end{prob}

\begin{proof}[न्यूनीकरणाने {\olref[nen-alt]{thm:nonenum-pownat}} ची सिद्धता]
न्यूनीकरणासाठी, $\Pow{\Nat}$ हा !!{enumerable} आहे आणि त्यामुळे त्याचे
$N_{1}$, $N_{2}$, $N_{3}$, \dots{} असे प्रगणन आहे, असे समजा.

$f \colon \Pow{\Nat} \to \Bin^\omega$ हे फलन पुढीलप्रमाणे परिभाषित
करा: $f(N)$ ही अशी चिन्हमाला~$s$ असू द्या की $s(n) = 1$ तेव्हा आणि
केवळ तेव्हाच, जेव्हा $n \in N$; अन्यथा $s(n) = 0$.
%READERNOTE{OLSIZ-010: गोठवलेल्या इंग्रजी व्याख्येत f(N) ला s_k असे म्हटले आहे, पण k बांधलेला किंवा N वरून ठरवलेला नाही. घटकनिहाय समीकरण N ची अभिलक्षण-चिन्हमाला एकमेव ठरवते; मराठी व्याख्येत एकच बद्ध निर्गत नाव s आणि त्याचे घटक s(n) सातत्याने वापरले आहेत.}

याने फलन स्पष्टपणे परिभाषित होते, कारण $N \subseteq \Nat$ असताना
कोणताही $n \in \Nat$ हा $N$ चा !!a{element} असतो किंवा नसतो.
उदाहरणार्थ, सम नैसर्गिक संख्यांचा संच
$2\Nat = \Setabs{2n}{n \in \Nat} = \{0,2, 4, 6, \dots\}$ हा
$1010101\dots$ या चिन्हमालेकडे; $\emptyset$ हा $0000\dots$ कडे; आणि
$\Nat$ हा $1111\dots$ कडे पाठवला जातो.

हे फलन !!{surjective} देखील आहे: $0$ आणि $1$ यांची प्रत्येक चिन्हमाला
नैसर्गिक संख्यांच्या कोणत्या तरी संचाशी अनुरूप असते; तो संच म्हणजे
चिन्हमालेत ज्या स्थानांवर~$1$ आहे त्यांना अनुरूप नैसर्गिक संख्या
सदस्य म्हणून असलेला संच. अधिक नेमकेपणाने, $s \in \Bin^\omega$ असेल,
तर पुढीलप्रमाणे $N \subseteq \Nat$ परिभाषित करा:
\[
N = \Setabs{n \in \Nat}{s(n) = 1}
\]
मग~$f$ ची व्याख्या पाहून $f(N) = s$ हे तपासता येते.

आता पुढील यादी विचारात घ्या:
\[
f(N_1), f(N_2), f(N_3), \dots
\]
$f$ हे !!{surjective} असल्यामुळे $\Bin^\omega$ चा प्रत्येक सदस्य
हा~$f$ चे कोणत्या तरी फलसाधकावरील मूल्य असला पाहिजे आणि म्हणून यादीत
आला पाहिजे. त्यामुळे ही यादी संपूर्ण~$\Bin^\omega$ चे प्रगणन करते.

म्हणून $\Pow{\Nat}$ हा !!{enumerable} असता, तर $\Bin^\omega$ हाही
!!{enumerable} असता. पण $\Bin^\omega$ हा !!{nonenumerable} आहे
(\olref[nen-alt]{thm:nonenum-bin-omega}). त्यामुळे $\Pow{\Nat}$ हा
!!{nonenumerable} आहे.
\end{proof}

%\begin{explain}
%न्यूनीकरण कोणत्या दिशेने केले जाते याबद्दल सहज गोंधळ होऊ शकतो.
%उदाहरणार्थ, $g \colon \Bin^\omega \to X$ असे !!a{surjective} फलन
%$X$ हा !!{nonenumerable} आहे, हे \emph{सिद्ध करत नाही}. ($g
%\colon \Bin^\omega \to \Bin$ हे $g(s) = s(1)$ ने परिभाषित केलेले,
%$0$ आणि $1$ यांच्या अनुक्रमाला त्याच्या पहिल्या !!{element} कडे
%पाठवणारे फलन विचारात घ्या. काही अनुक्रम $0$ ने आणि काही $1$ ने सुरू
%होत असल्यामुळे ते आच्छादक आहे. पण $\Bin$ हा सांत आहे.) फलन~$f$ हे
%आच्छादक असलेच पाहिजे, हेही लक्षात घ्या; अन्यथा युक्तिवाद लागू होत नाही:
%$f(x_1)$, $f(x_2)$, \dots{} मध्ये~$Y$ चे सर्व !!{element}s असतील
%याची खात्री उरत नाही. उदाहरणार्थ, पुढीलप्रमाणे परिभाषित केलेले
%$h\colon \Nat \to
%\Bin^\omega$:
%\[
%h(n) = \underbrace{000\dots0}_{\text{$n$ वेळा $0$}}
%\]
%हे फलन आहे, पण $\Nat$ हा !!{enumerable} आहे.
%\end{explain}

\begin{prob}\label{sfr:siz:red:prob:nat-nat}
  $f\colon \Nat \to \Nat$ अशा सर्व फलनांचा संच~$X$ हा
  न्यूनीकरणाच्या युक्तिवादाने !!{nonenumerable} आहे, हे दाखवा.
  (सूचना: $X$ पासून~$\Bin^\omega$ कडे जाणारे आच्छादक फलन द्या.)
\end{prob}

\begin{prob}
नैसर्गिक संख्यांच्या जोड्यांच्या \emph{सर्व संचांचा} संच, म्हणजे
$\Pow{\Nat \times \Nat}$, न्यूनीकरणाच्या युक्तिवादाने
!!{nonenumerable} आहे, हे दाखवा.
\end{prob}

\begin{prob}
नैसर्गिक संख्यांच्या अनंत अनुक्रमांचा संच $\Nat^\omega$ हा
न्यूनीकरणाच्या युक्तिवादाने !!{nonenumerable} आहे, हे दाखवा.
\end{prob}

%\begin{prob}
%$P$ हा $\Nat$ पासून $\{0\}$ या संचाकडे जाणाऱ्या फलनांचा संच असू द्या;
%आणि $Q$ हा धन पूर्णांकांच्या संचापासून $\{0\}$ या संचाकडे जाणाऱ्या
%\emph{अंशतः} फलनांचा संच असू द्या. $P$ हा !!{enumerable} आहे आणि
%$Q$ नाही, हे दाखवा. (सूचना: $\Bin^\omega$ चे प्रगणन करण्याची समस्या
%~$Q$ चे प्रगणन करण्याच्या समस्येकडे न्यूनीत करा.)
%\end{prob}

\begin{prob}
$S$ हा $\Nat$ पासून $\{0,1\}$ या संचाकडे जाणारी सर्व
!!{surjection}s असलेला संच असू द्या; म्हणजे $S$ मध्ये
$f \colon \Nat \to \Bin$ अशी सर्व !!{surjection}s आहेत. $S$ हा
!!{nonenumerable} आहे, हे दाखवा.
\end{prob}

\begin{prob}
सर्व वास्तव संख्यांचा संच~$\Real$ हा !!{nonenumerable} आहे, हे दाखवा.
\end{prob}

\end{document}
